% !TEX TS-program = xelatex
\documentclass[12pt,letterpaper]{article}

\input{../../../../../.house-style/preamble.tex}
\setlength{\headheight}{26pt}
\setlength{\emergencystretch}{2em}

\hypersetup{
  pdftitle={Why some interpretable deviations count as ungrammatical: grammatical-choice profiles and public-update consequences},
  pdfkeywords={grammaticality, form-meaning pairing, grammatical choice, operator profile, public update, repair}
}

\title{Why some interpretable deviations count as ungrammatical:\\grammatical-choice profiles and public-update consequences}
\author{Brett Reynolds \orcidlink{0000-0003-0073-7195}\\Humber Polytechnic \& University of Toronto}
\date{\today}

\begin{document}
\maketitle

\begin{abstract}
Why can an intelligible deviation be classified as ungrammatical in one case,
but as odd meaning, bad wording, accent, style, or social conduct in another?
I propose that such classifications fallibly track the relation that failed in
a community's inventory of conventional form--meaning pairings. A grammatical
choice isn't merely any conventional system: it recurs at a constructional
choice site, relates expression to a language-specific value or function, and
is constrained by a community licensing pattern. Strict paradigms are a strong
case, but construction-sized and partly lexicalized relations create graded and
moving edges. A token can lack a viable analysis, combine incompatible values,
leave an obligatory contrast unsupplied, select the wrong paradigm value, or
use an unlicensed realization. By contrast, a poor lexical, stylistic, or social
choice may preserve the relevant form--meaning pairings and instead mis-select
content, persona, situation, or conduct. Production, substitution, omission,
semantic, and distributional evidence can place these relations without using
the fault labels to be predicted. The primary claim is prospective:
stronger independently assigned \term{grammatical-choice profiles} should
predict more consistent grammar-directed attribution in held-out cases. Within
that architecture, a graded public-update profile identifies one consequential
subset. Building on Heltoft and N{\o}rg{\aa}rd-S{\o}rensen, I predict that its
strength increases the difference in update-directed repair between wrong-value
and value-preserving-realization errors. Abrus\'an's semantic-repairability
account, Boye and Harder's discursive secondary status, processing, schooling,
and sanction remain live alternatives. Either projection can fail without
turning the other into a definitional success.
\end{abstract}

\smallskip
\noindent \textbf{Keywords:} grammaticality; form--meaning pairing; grammatical choice; operator profile; public update; repair

\smallskip
\noindent\footnotesize\textit{Version note.} An earlier version posted as
LingBuzz/009706 proposed a broader coordination account of selective
grammaticality. This article replaces its categorical operator-stratum
taxonomy with a graded profile over independently identified form--value
assembly relations and
recasts repair as a downstream test of a broader account of selective
grammatical classification.\normalsize

\aidisclosure{Claude Opus 4.5 and Claude Fable; a preview model with no
disclosed developer or underlying-model identity, accessed through OpenCode Zen
in August 2026; ChatGPT 5.2 Pro; and GPT-5.5}


\section{The selective-classification problem}

The initiating example looks too small to carry a theory. In ordinary connected
Reference French, \olang{hiver} `winter' begins with \term{h muet}, so the
singular definite article undergoes obligatory elision: \olang{l'hiver}, not
*\olang{le hiver} \citep{boersma2007,simonenkocarlier2020}. Yet the deviant form
remains readily interpretable. \olang{Le} still expresses masculine singular
definiteness; what fails is its conventional morphophonological realization.
The deviation is neither a wrong referent nor a wrong definiteness value. Nor is
it merely any French-accent difference. The realization relation belongs to the
article construction, and that location makes the example a candidate
grammatical form error.

This case generated the paper's question. Why do some interpretable departures
count as grammar, while others count as semantic anomaly, poor wording,
pronunciation, style, or social conduct? The contrast can't be reduced to
meaninglessness: *\olang{le hiver}, *\mention{goed}, and many agreement errors
are transparent. Nor can it be reduced to expressive substance. Tone,
intonation, word order, gesture, and segmental morphology can all realize
grammatical contrasts, while segmental differences can also mark accent or
style. An asterisk records a classification; it doesn't yet explain that
classification.

Two questions are often run together. \term{Source attribution} asks what
caused a negative response: grammatical knowledge, processing load, frequency,
schooling, social evaluation, or something else. \term{Target classification}
asks what the respondent treats as having failed: a construction, a grammatical
value, its realization, lexical content, style, stance, or the social act. The
acceptability literature's source-ambiguity problem concerns the first question
\citep{hofmeisteretal2013source,leivadawestergaard2020}. This paper asks
how the second can be predicted without reading the answer from the response it
is meant to explain.

The proposed answer begins from a familiar but insufficient observation:
grammar consists centrally of conventional pairings between expression and
content or function, at grains ranging from words to abstract constructions
\citep{goldberg1995constructions,sag2012sign}. Phonemic contrasts ordinarily
keep expressions apart rather than pairing each contrast with a meaning;
orthography ordinarily maps graphic forms to linguistic expressions. But the
pairing criterion alone can't separate grammar from the lexicon or from social
meaning. A word and an enregistered style also pair form with value.

The decisive question is relational: \emph{which relation does the
departure fail?} Selecting an inapposite but otherwise licensed word preserves
that word's form--meaning pairing; a register mismatch can preserve both the
proposition and the social value conventionally indexed by the form. Trouble is
then directed at chosen content, situational fit, persona, or conduct. A
grammatical departure instead compromises the conventional assembly of
form--meaning relations: an analysis is unavailable, grammatical instructions
conflict, an obligatory contrast is missing, a construction is excluded, a
paradigm value is mis-selected, or an intended value receives an unlicensed
realization. Lexical and social material can enter grammar when recruited
into such a constructional relation; morphology, tone, or word order doesn't
count as grammar merely by virtue of its substance.

This answer would be circular if \enquote{grammatical relation} were assigned
from the same asterisks and corrections it later explained. I propose
a graded \term{grammatical-choice profile}, assigned from independent evidence:
semiotic and constructional assembly, contrastive choice architecture, and
conditioned community licensing. These properties are neither a classical
checklist nor a synonym for \enquote{system}.
They order candidate relations for a prospective claim: higher independently
assigned profiles should predict more consistent grammar-directed fault
attribution in held-out cases. A stable failure defeats or splits that claim;
it can't be repaired by moving the category edge after the labels are seen.

Strict paradigms in the sense developed by N{\o}rg{\aa}rd-S{\o}rensen and
Heltoft occupy one strong region of this profile space. Constructional licensing
supplies another, often layered, route. Within either route, a token may fail at
the level of coverage, compatibility, saturation, community licensing, value,
or realization. This imports the main separations of OVMG without importing its
population dynamics: OVMG distinguishes the status of a candidate from the
evidence available to a speaker and from the resulting experience or label
\citep{reynolds2026ovmg}. This paper asks what kind of linguistic target
that label purports to identify.

Public-update function then contributes a narrower projection rather than a
definition of grammar. A graded \term{public-update profile} measures how
strongly a paradigm's values conventionally configure move type, commitment or
source, discourse continuity, and response relevance. A wrong
switch-reference value can change which participant is carried forward; a
malformed realization of the same intended value need not. The
\term{operator-profile hypothesis} predicts that this value--realization
difference in update-directed repair grows with a profile assigned before
held-out responses are coded. *\olang{le hiver} matters precisely because it
shows why this operator-specific claim can't define grammaticality: the
definite value survives even though its realization is excluded.

The paper thus proposes three separable projections. The primary one concerns
selective fault attribution across grammatical, lexical, phonological,
orthographic, stylistic, pragmatic, and social targets. The second concerns
variable and moving attribution at intermediate or changing edges. The third
concerns the downstream consequences of wrong values within a subset of
grammar. Repeated
opportunities and competitor strength can explain why a transparent form
becomes sharply excluded \citep{reynolds2026ovmg}; they don't by themselves
decide which relation failed or where subsequent trouble will be directed.

Section~\ref{sec:prior-constraints} positions that test against the closest
partial answers. Sections~\ref{sec:architecture} and
\ref{sec:profiles} develop the choice profile and the operator-specific
coordinate. Section~\ref{sec:comparisons} develops worked comparisons.
The final sections compare rival predictors and state held-out designs in which
each projection can fail.

\section{Prior constraints and live alternatives}
\label{sec:prior-constraints}

Abrus\'an's programme is the closest precedent in problem shape.
\textcite[329--331, 344--348]{abrusan2019} asks why failed interpretation can
produce semantic anomaly, pragmatic infelicity, or ungrammaticality.
\textcite{abrusanetal2021shift} develop an answer in terms of semantic
repairability: local conflicts involving lower-level, context-shiftable types
can leave a recoverable predication, whereas conflicts involving constitutive
high-level types can be nonlocal and resist meaning-preserving shift. The type
architecture and shift predictions are independently motivated, yielding a
real explanation of one classification boundary. Its native explanandum is
semantic anomaly and the semantically or pragmatically analyzed subset of
ungrammaticality. It doesn't purport to classify article elision, suppletive
allomorphy, constructional exclusion, accent, or social sanction.

The source-ambiguity literature supplies a complementary constraint. A low
rating doesn't reveal whether grammar, processing, interpretation, or
contextualization produced it \citep{hofmeisteretal2013source}; grammaticality
and acceptability can dissociate in both directions
\citep{leivadawestergaard2020}. In functional terms, grammaticality concerns
fit to relatively stable conventional knowledge rather than acceptability in
one act of use, but corpus distributions, processing measures, and judgments
are all indirect evidence for that knowledge
\citep[12--20]{poulsen2012grammaticality}. No single response may be
used both to locate a fault and to explain why the fault was located there.

Coseriu provides a deeper historical precedent for differentiated fault
domains. His system--norm--speech distinction separates functional
possibilities, customary social realization, and individual production, while
later presentations distinguish congruence, idiomatic correctness, and
situational appropriateness \citep{coseriu1952,pardollibrer2021,sotovergara2022}.
This paper doesn't rediscover differentiated evaluation. It asks which
independently measurable properties predict a particular allocation in later
behaviour.

Construction Grammar supplies the broad semiotic ontology. Words, idioms,
partly filled patterns, and abstract constructions are conventional
form--meaning pairings rather than outputs of a grammar divided cleanly from a
lexicon \citep{goldberg1995constructions,sag2012sign}. That inheritance rules
out an easy substance boundary, but it leaves the present question open. A
licensed word can be badly chosen without its pairing failing; a construction
can be unlicensed even when its content is transparent. The distinction to be
explained concerns relations and fault targets inside a structured
inventory, not between meaningful words and meaningless syntax.

Heltoft's work is the closest positive scaffold. His analyses distinguish
grammaticality from comprehensibility and treat mood, word order, particles,
and other paradigms as resources for illocutionary and interlocutionary
organization
\citep{heltoft1981intersubjectivity,heltoft2005modality,hansenheltoft2011gds,jakobsenheltoft2012paradigm}.
\textcite{norgardsorensenheltoft2015} generalize the paradigm architecture:
bounded alternatives occupy a recurring domain and contrast within a
language-specific semantic frame. Their architecture reaches beyond
morphology, and their Danish V2 analysis already joins paradigm structure to
assertivity. This paper adopts that scaffold. It doesn't claim that either
paradigms or interactional grammar are new. Their additional marked--unmarked
asymmetry anchors the strict paradigmatic core considered below. It neither
defines all grammar nor warrants an inference from paradigmhood to categorical
judgment. Less bounded, construction-sized, or item-restricted relations can
share some of the profile without satisfying the full paradigm definition.

Boye and Harder are the principled alternative, not a minor comparison. They
treat closed paradigms and obligatoriness as incidental to grammatical status;
conventional ancillarity and discursive secondary status are constitutive, with
focusability and addressability as diagnostics \citep{boyeharder2012}. Their
account directly competes with the grammatical-choice profile as an account of
status. Its diagnostics can also compete with the public-update profile in the
narrower repair study. If discursive secondary status predicts held-out fault
attribution or subsumes the proposed repair interaction, the corresponding new
profile hasn't earned an additional explanatory role.

\begin{table}[ht]
\centering
\scriptsize
\setlength{\tabcolsep}{3pt}
\begin{tabular}{@{}p{0.20\linewidth}p{0.25\linewidth}p{0.25\linewidth}p{0.22\linewidth}@{}}
\toprule
\textbf{Programme} & \textbf{Native explanandum} & \textbf{Independent predictor} & \textbf{Residual question here} \\
\midrule
Abrus\'an and colleagues & semantic anomaly, pragmatic infelicity, and semantically driven ungrammaticality & type level, locality, and meaning shift & form, construction, style, and social-act targets \\
Coseriu & system, normal realization, speech, and evaluative domains & functional opposition and norm-relative realization & held-out allocation and gradience of token faults \\
Construction Grammar & learned inventory across lexicon and syntax & conventional form--meaning pairings at multiple grains & why one licensed pairing is badly chosen while another pairing relation fails \\
Heltoft and N{\o}rg{\aa}rd-S{\o}rensen & language-specific grammatical paradigms and their content & domain, frame, bounded alternatives, obligatory choice, markedness & graded extension beyond strict paradigms and held-out attribution \\
Boye and Harder & grammatical status and grammaticalization & discursive secondary status & comparison with choice profile and repair locus \\
Source ambiguity and functional acceptability & causes of ratings and evidence for stable conventional knowledge & triangulated grammar, processing, and context measures & target classification rather than response cause \\
\bottomrule
\end{tabular}
\caption{The closest partial answers. The final column states a difference in
scope or design, not a defect in the predecessor.}
\label{tab:predecessors}
\end{table}

Interactional grammar and repair research shows that grammar is a resource for
sequential action and supplies coding precedents for downstream trouble and correction
\citep{foxetal1996repair,muntiglventola2010,dingemanseblythedirksmeyer2014}.
Update semantics and clausal-operator accounts supply the dimensions of move,
commitment, and response potential used below
\citep{stalnaker2002,farkasbruce2010,murray2014update,hengeveldmackenzie2008fdg,bentleyetal2023rrg}.
Neither literature already contains the two frozen comparisons proposed here:
from a graded choice profile to fault-domain attribution, and from the
public-update coordinate to the value--realization repair difference. That
residual claim is architectural and methodological, not the discovery that
grammar is semiotic, paradigmatic, or interactionally consequential.

\section{Cross-level architecture}
\label{sec:architecture}

The neutral bearer of the proposal is a candidate \term{form--value assembly
relation} in a conditioned community repertoire. Calling the bearer a
\enquote{grammatical system} would assume the classification to be explained;
calling it merely a \enquote{system} would include vowel inventories,
orthographies, politeness routines, and practices of apologizing. A candidate
relation instead specifies an expression, a semantic or relational value, the
larger construction in which they combine, and the community and situations in
which the dependency is claimed to hold.

Its \term{grammatical-choice profile} records three blocks:

\begin{enumerate}
  \item \term{Semiotic and constructional assembly}: conventional coupling
    between expression and semantic, relational, or structural value, recurring
    at a specifiable site in larger expressions.
  \item \term{Choice architecture}: alternatives organized in a
    language-specific domain and frame, with their boundedness and conditioned
    requirements on selection, omission, or realization recorded.
  \item \term{Conditioned licensing}: the relation's position and generality in
    a multi-grain constructional chain, and the concentration of its use across
    the declared community, situations, medium, and period.
\end{enumerate}

Public-update contribution is measured alongside this structural profile in
\S\ref{sec:profiles}; it isn't smuggled into every grammatical classification.
No block is a disguised judgment score. Positive production distributions,
substitution and omission contrasts, semantic interpretation, productivity,
and opportunity-normalized licensing evidence place the candidate before the
fault labels, corrections, or repairs to be predicted are inspected. A
descriptive grammar contributes only through evidence of those relations, not
by lending the study its prior use of \enquote{grammatical}. The theoretical
paper proposes no universal scalar; concrete studies have to preregister ordinal or
pairwise comparisons, or a composite and its weighting rule.

The profile is deliberately graded. A strict paradigm has a recurring domain,
a language-specific frame, a bounded set of alternatives, obligatory value
selection in the domain, and marked--unmarked organization
\citep[262--263]{norgardsorensenheltoft2015}. It anchors a strong
region of the space. A productive constructional dependency may be less
bounded; an item-specific complement restriction may sit deep in a grammatical
licensing chain while resembling lexical collocation at its narrowest link.
These needn't be equal cases, and neither has to be excluded by a
classical definition before its projections are tested. Hansen and Heltoft's
own grammar permits prototype and radial organization
\citep[vol.~1, 39--41]{hansenheltoft2011gds}.

Three kinds of gradience need to remain distinct. The worldly dependency can have
an intermediate profile; evidence for its placement can be uncertain; and a
speaker or institution can issue a sharp or hesitant label. A categorical
classroom correction doesn't prove a strong grammatical profile, and uncertain
evidence doesn't prove a fuzzy system. Profiles and category edges can also
move across communities, activities, media, and time. Those changes are part of
the claim's scope, not exceptions to be hidden after a miss.

OVMG supplies a second decomposition at the token level. Once a candidate
analysis and its conditioning situation are fixed, \term{coverage} asks whether
an analysis exists; \term{compatibility} whether its grammatical instructions
can jointly hold; \term{saturation} whether required contrasts have values; and
\term{licensing} whether the community supports the construction and
realization. Population status, evidence available to a speaker or analyst, and
the resulting subjective read-out remain separate \citep{reynolds2026ovmg}.
This paper uses those distinctions to locate a possible fault without
reproducing OVMG's dynamics.

\begin{table}[ht]
\centering
\scriptsize
\setlength{\tabcolsep}{3pt}
\begin{tabular}{@{}p{0.19\linewidth}p{0.24\linewidth}p{0.31\linewidth}p{0.18\linewidth}@{}}
\toprule
\textbf{Question} & \textbf{Bearer} & \textbf{Possible failure} & \textbf{Example} \\
\midrule
Conditioning & Candidate plus situation and norm centre & wrong repertoire or context has been assumed & formal/informal contrast \\
Coverage & Proposed constructional analysis & no viable assembly of the supplied material & *\mention{Can the have running?} \\
Compatibility & Values in a covered analysis & independently supplied instructions conflict & *\mention{I've finished it yesterday} \\
Saturation & Obligatory contrast in a covered analysis & required value remains unsupplied & omitted obligatory marker \\
Licensing & Construction or realization in a repertoire & transparent candidate is conventionally excluded & *\mention{depend of}; *\olang{le hiver} \\
Value selection & Token of an identified paradigm & licensed competitor selects the wrong value here & same vs.\ different subject \\
Non-assembly target & Licensed sign or social act & content, spelling, stance, persona, situation, or conduct is at issue & register mismatch \\
\bottomrule
\end{tabular}
\caption{Cross-level diagnostic questions. Value and realization specify loci
inside the broader compatibility and licensing questions. A response can target
more than one row, and its cause is separate.}
\label{tab:strata}
\end{table}

The primary projective claim, \term{selective fault attribution}, now has a
non-trivial form. For candidate dependencies in a declared community,
situation, medium, and period, a stronger profile assigned without attribution
evidence should predict more consistent grammar-directed classification of
held-out interpretable deviations. Frequency, processing load, semantic
repairability, discursive secondary status, schooling, and sanction are rival
predictors. Before a concrete study, the minimum useful improvement and
uncertainty rule has to be declared. If the profile adds no held-out reach, it's
demoted or split; the misses may not be excluded by redescribing the relation
after the responses are seen.

The account predicts attribution; respondents can still classify incorrectly.
An institutional grammar label can target a weakly
profiled shibboleth, while an unprimed speaker can treat a strong but changing
dependency as lexical or stylistic. The proposal predicts a robust core,
layered interfaces, and moving edges rather than exceptionless agreement.

This makes *\olang{le hiver} informative. Its definite, masculine, singular values
remain recoverable. The failure is the unelided realization inside an ordinary
French determiner phrase, not the selection of a different value. The analysis
is deliberately neutral between deriving \olang{l'} by vowel deletion and
treating it as an allomorph: both locate the error in morphophonological
realization rather than concord \citep{boersma2007}. The example establishes
neither a wrong value, the containing system's paradigm status, nor its
public-update profile. It functions instead as an interface test. Spoken schwa retention,
written non-elision, learner speech, and a deliberate pause can make different
dependencies and norm centres salient, so the account permits mixed
phonological, orthographic, learner, or grammatical-form attribution. In OVMG
terms, the intended definite assembly is covered, compatible, and saturated;
the disputed relation is realization licensing in the specified environment.

Repeated opportunities help explain why such a transparent failure can be
sharp. Once the domain and alternatives are independently specified, a small
competitor set allows evidence to concentrate quickly on the licensed option.
OVMG develops that licensing and preemption account
\citep{reynolds2026ovmg}. Information theory can quantify the resulting
uncertainty reduction, but it can't discover the grammatical alternatives or
decide which fault domain a community recognizes. The linguistic codebook is an
input to that calculation, not its output \citep{shannon1948,coverthomas2006}.

For the narrower operator-profile comparison below, study eligibility is
categorical even though grammaticality isn't. A candidate paradigm has to have
bounded alternatives, a recurring domain, a shared frame, and obligatory value
selection; markedness is recorded, with a strict five-criterion sensitivity
analysis. This gate prevents an open set from entering a value--realization
experiment merely because some members matter to interaction. It's a design
rule for that comparison, not a definition of grammar.

\section{Public-update profiles}
\label{sec:profiles}

Fault location doesn't yet predict what a hearer will do with the trouble. An
identified paradigm receives a graded profile for the extent to which
its values conventionally configure four aspects of public update: move type,
commitment holder or source, discourse-participant or referent continuity, and
response relevance \citep{stalnaker2002,farkasbruce2010,murray2014update}. The
components form a vector, not four additional membership conditions. Clause
type, evidential anchoring, and switch-reference are strong candidates;
neither thematic-role assignment nor formal obligatoriness is sufficient on its own.

The profile belongs to the paradigm across its ordinary occasions of use, not
to the information carried by every token. An agreement marker can be redundant
in one clause and decisive in another. Licensed update type also differs from
truth or success: a false assertion is still an assertion, while a wrong
same-subject value can redirect reference tracking even when the speaker's
intention is recoverable.

Within identified paradigms, a \term{wrong-value} error changes the grammatical
choice. A \term{value-preserving-realization} error retains that choice but uses
an excluded exponent or surface form. The \term{operator-profile hypothesis}
states that, for paradigms in a specified community and activity range that
independently support both error types, a stronger frozen profile predicts a
larger held-out difference in update-directed repair between them. The
within-paradigm contrast checks the manipulation; the cross-paradigm interaction
is the explanatory claim.

The bridge is a conditional claim about sufficient repair. When trouble is
detected and participants pursue mutual understanding, a response can target
the lowest locus whose correction settles the public record. A wrong value in a
high-profile paradigm can leave move type, commitment, reference continuity, or
response relevance unsettled; correcting surface form alone may be
insufficient. When the intended value is independently recoverable and only its
realization is excluded, supplying the licensed form can settle the trouble
without renegotiating the update. This doesn't make overt repair inevitable.
Accommodation, high recoverability, non-response, mishearing, and person-directed
sanction are counterpressures and remain in the outcome space.

The information-theoretic intuition concerns a code's organization, not merely
token surprisal. Recurrent constrained choices create redundancy: a hearer can
recover the intended value precisely because the licensed competitors make the
deviation easy to detect. Recoverability doesn't show that the
violated relation carries no grammatical information. Public-update profiling
asks a different counterfactual question: how much would choosing another
licensed value change the move, commitment, continuity, or relevant response
space? Entropy or divergence can summarize that change only after linguistic
analysis supplies the alternatives \citep{shannon1948,coverthomas2006}.

I retain \term{operator stratum} only as shorthand for the higher-profile region
of this ordering. It introduces no threshold or second membership category.
Repair and comprehension are outcomes, never profile criteria. A high-profile
paradigm may be accommodated, ignored, sanctioned, or corrected only in form;
those responses count against the hypothesis rather than licensing a post-hoc
redescription.

Opportunity remains distinct. It affects how concentrated the evidence for a
licensed realization becomes, not whether a contrast is a paradigm or how
strongly its values configure public update. Likewise, historical recruitment
is a separate question. Grammaticalization research documents the recruitment
of material into grammatical and discourse functions, often through
paradigmatic reorganization
\citep{hoppertraugott2003,lehmann2015grammaticalization,waltereit2011discourse,degandeversvermeul2015,jakobsenheltoft2012paradigm}.
It doesn't establish that public-update work preferentially creates paradigms,
and diachronic conjecture can't rescue a failed synchronic prediction.

Expressive substance is also independent. Morphology, word order, particles,
tone, intonation, and nonmanual marking can all realize values in identified
paradigms. Phonemic systems can be tightly organized while serving primarily to
distinguish expressions; open interjectional resources can strongly configure
uptake without forming a bounded paradigm \citep{ameka1992}. Paradigm structure,
public-update profile, and expressive substance answer different questions.

\section{Worked comparisons and boundary anchors}
\label{sec:comparisons}

The architecture needs examples that keep value and realization apart. It also
needs boundary anchors that expose rather than stipulate the proposed category
edge. Table~\ref{tab:boundary-cases} records candidate analyses and attributions;
none is an observed result.

\begin{table}[ht]
\centering
\scriptsize
\setlength{\tabcolsep}{3pt}
\begin{tabular}{@{}p{0.19\linewidth}p{0.25\linewidth}p{0.25\linewidth}p{0.23\linewidth}@{}}
\toprule
\textbf{Case} & \textbf{Failed dependency} & \textbf{Profile expectation} & \textbf{Candidate attribution} \\
\midrule
Wrong paradigm value & value selection in a recurring domain & strong where strict paradigm criteria hold & grammar/value; possibly update \\
*\olang{le hiver} & conditioned article realization & construction-embedded but phonology- and medium-sensitive & grammatical form, phonology, or orthography \\
*\mention{goed} & past realization for one lexeme & inflectional choice plus item-specific licensing & grammatical form or word form \\
*\mention{depend of} $+$ \textsc{obj} & item-restricted link in a complement chain & grammatical at higher links, collocational at narrowest link & construction or wording \\
Register or accent mismatch & situational fit of a licensed indexical value & form--value coupling without necessary assembly failure & style, persona, or social norm \\
Misspelling & graphic realization of a linguistic expression & representation relation normally preserves linguistic assembly & orthography \\
\bottomrule
\end{tabular}
\caption{Boundary anchors for selective attribution. Profile expectations have to
be fixed independently; the final column is a prediction, not a classification
criterion.}
\label{tab:boundary-cases}
\end{table}

The French example remains the motivating realization case, but it isn't a
concord case and supplies no wrong-value counterpart. *\mention{Goed} is
parallel at a coarser, lexeme-specific grain: the past value survives while the
established exponent \mention{went} displaces the productive default.
*\mention{Depend of} is layered. Complementation selects a function; that
function constrains a phrasal category; the PP option introduces an
item-restricted preposition choice. The higher links have strong schematic and
assembly profiles, whereas the final \mention{depend}--\mention{on} restriction
approaches lexical collocation. Mixed construction-versus-wording attribution
would support the layered analysis rather than count as noise.
Sharp correction of any of the three would by itself establish neither the
source profile nor a high operator profile.

\subsection{Amele finite tense}

Amele supplies the strongest source-grounded template presently available. Its
finite verb paradigm distinguishes present, today past, and yesterday past; the
grammar also states temporal co-occurrence restrictions and gives the
morphophonological derivations of the forms
\citep[226--232, 380--385]{roberts1987}. Holding the verb `come' and first-person
singular subject constant yields the comparison in
Table~\ref{tab:amele-tense}.

\begin{table}[ht]
\centering
\small
\begin{tabular}{@{}p{0.22\linewidth}p{0.30\linewidth}p{0.38\linewidth}@{}}
\toprule
\textbf{Condition} & \textbf{Amele form} & \textbf{Analysis} \\
\midrule
Licensed target & \olang{Ija hu-g-a} & `I came today/have come' \\
Wrong value & \olang{Ija hu-g-an} & licensed yesterday-past form in an unambiguously earlier-today context \\
Preserved value, wrong realization & *\olang{Ija ho-ig-a} & today-past and 1sg material retained, but obligatory vowel coalescence fails \\
\bottomrule
\end{tabular}
\caption{A proposed Amele finite-tense comparison. The final form is derived
from the published morphophonology and requires speaker validation.}
\label{tab:amele-tense}
\end{table}

The value error substitutes a licensed competitor from the same tense
paradigm. The realization error retains the today-past suffix \olang{-a} and
first-person singular material but fails the derivation
\olang{/ho+ig+a/ $\rightarrow$ [huga]}. Roberts independently stars a comparable
nonapplication or misordering outcome elsewhere in the finite paradigm, which
supports the analysis without supplying a judgment for this exact item
\citep[380--385]{roberts1987}. The result is a theoretically clean
template, not a ready-made experimental result. A future use has to validate the
constructed form, equate perceptual distance, and avoid the flexible nighttime
boundary between today and yesterday past.

Nothing in the grammatical label assigns the tense paradigm a high profile.
Study B Stage~1b has to determine how strongly its values alter temporal anchoring,
commitment, or response relevance. The worked case establishes that value and
realization can be manipulated separately.

\subsection{Amele switch-reference}

The sequential switch-reference paradigm provides a second template. Amele
distinguishes same-subject (SS) and different-subject (DS) linkage on a medial
verb, with separate person--number sets and vowel harmony in the DS marker
\citep[239--240, 298--300, 371]{roberts1987}. For a second-person subject
followed by a third-person subject, the comparison is:

\begin{table}[ht]
\centering
\small
\begin{tabular}{@{}p{0.22\linewidth}p{0.36\linewidth}p{0.32\linewidth}@{}}
\toprule
\textbf{Condition} & \textbf{Amele form} & \textbf{Analysis} \\
\midrule
Licensed target & \olang{Hina ho-co-m sab je-i-a} & `You came and he ate the food' (DS) \\
Wrong value & *\olang{Hina h-u-me-g sab je-i-a} & licensed 2sg SS form with a different following subject \\
Preserved value, wrong realization & *\olang{Hina ho-ce-m sab je-i-a} & DS and 2sg material retained, but the DS vowel is disharmonic after \olang{ho-} \\
\bottomrule
\end{tabular}
\caption{A proposed Amele switch-reference comparison. Both starred sentences
require validation in the intended clause-chain context.}
\label{tab:amele-sr}
\end{table}

Here the wrong-value token changes the tracking instruction, while the
disharmonic token preserves DS. The grammar directly supports the licensed SS
and DS forms and the harmony relation, but it doesn't print the exact
*\olang{ho-ce-m} stimulus. SS/DS choice is also sensitive to agentivity,
clause level, time, place, world, and pragmatic deixis. Those factors have to be
frozen rather than treated as noise. Because SS and DS use different
person--number material, the two error conditions also differ in more than one
segment; perceptual salience belongs in the design.

Switch-reference is a plausible high-continuity profile, but that expectation
isn't a score. A separate profiling partition has to establish whether changing
SS to DS alters reference continuity and downstream response options. The
worked comparison shows what an eligible case would look like and where its
remaining evidential risks lie.

\subsection{Neighboring targets and moving edges}

Gender and noun-class concord stress-test the paradigm-level analysis. A token
such as *\olang{le table} may leave the intended referent obvious while violating
obligatory French concord. Local redundancy doesn't show that the containing
paradigm lacks a reference-tracking role, but neither does categorical correction
prove that it has one \citep{corbett1991,mahowaldetal2023}. Independent profiling
has to compare contexts in which lexical semantics already fixes the referent
with contexts in which concord changes participant tracking.

Accent, register, and Japanese polite/plain style supply low-profile or
non-paradigm comparisons only after their conditioned repertoires are specified
\citep{agha2007,eckert2012,cook2008speechstyle,ide2005}. Their deviations can be
highly accountable while targeting stance, persona, or institutional norms.
The form can express exactly the social value conventionally associated with it
and still be wrong for the occasion. In that case the pairing has succeeded;
situational fit is the target. But an honorific or T/V contrast that acquires a
recurring domain, bounded alternatives, and conditioned selection can move
toward the grammatical core. The theory has to allow that result rather than
assigning indexical meaning outside grammar by stipulation.

Phonology and orthography sharpen the same point. A phonemic contrast ordinarily
serves sign discrimination rather than pairing each segment with a meaning. A
spelling system ordinarily maps graphic forms to linguistic expressions. Their
departures can leave the relevant linguistic form--value assembly
intact. Morphophonology and grammatical tone cross the boundary when a
phonological contrast realizes a constructional or paradigm value. *\olang{le
hiver} is informative because respondents may locate the same surface departure
at that interface differently across speech, writing, learner correction, and
deliberate pauses.

The lexicon contains semiotic pairings. Words are form--meaning pairings, but
choosing an ill-fitting licensed word normally preserves its pairing and the
larger assembly. Selectional restrictions, valence frames, and item-specific
complement patterns create intermediate cases when lexical material constrains
how a construction is assembled. *\mention{Depend of} is useful only
if the licensing chain is specified in advance; calling it simply lexical or
grammatical would throw away the prediction.

These interfaces motivate the second broad projection, \term{moving-edge
attribution}. Independently measured changes in recurrence, constructional
embedding, boundedness, conditioned selection, or population concentration
should predict movement and increased variation among grammatical, lexical, and
stylistic attributions. The source profile may change historically or differ
across activities; those comparisons require separately declared populations
and periods. Schooling is a rival cause: an institution can preserve a grammar
label after spontaneous production and licensing have shifted. Unprimed
attribution and explicit schoolroom classification have to be measured
separately.

Interjections and slurs make the converse point: a resource can reshape uptake
or social action without entering a bounded grammatical paradigm
\citep{ameka1992}. These cases prevent public consequence from becoming a
shortcut to either grammatical-choice or public-update profile.

\section{Incremental predictor comparison}
\label{sec:predictor-comparison}

The nearest predictors have to be compared within their native scope rather
than forced into a winner-take-all theory of language. Table~\ref{tab:predictors}
states what each can directly predict in the proposed study.

\begin{table}[ht]
\centering
\scriptsize
\setlength{\tabcolsep}{3pt}
\begin{tabular}{@{}p{0.19\linewidth}p{0.23\linewidth}p{0.26\linewidth}p{0.24\linewidth}@{}}
\toprule
\textbf{Predictor} & \textbf{Applicability gate} & \textbf{Direct prediction} & \textbf{Probative result} \\
\midrule
Abrus\'an: type level, locality, shiftability & semantic or pragmatic analysis of the anomaly & whether local meaning repair preserves a coherent predication & explains semantic targets and absorbs apparent profile effects \\
Boye--Harder: discursive secondary status & expression stands in the relevant syntagmatic relation & grammatical status and associated focusability/addressability pattern & secondary status predicts attribution beyond the choice profile or repair beyond paradigm structure \\
Grammatical-choice profile & independently profiled form--value assembly relation & consistency and domain of grammar-directed attribution & incremental held-out allocation beyond frequency, processing, schooling, and sanction \\
Public-update profile & independently identified paradigm with both error types & value/realization difference in update-directed repair grows with profile & incremental interaction after status rivals and licensing variables are represented \\
\bottomrule
\end{tabular}
\caption{Scope-fair predictor comparison. An N/A case is outside a predictor's
native explanandum, not evidence against it.}
\label{tab:predictors}
\end{table}

Abrus\'an's account may fully explain items whose apparent grammaticality is a
semantic composition failure. Boye and Harder may better identify grammatical
status generally. The grammatical-choice profile earns its broader role only
if its independently measured architecture predicts held-out target allocation.
The operator profile earns its narrower role only if it then adds reach for the
consequences of value errors. A broad interpretive-consequence measure also has to
be allowed to subsume it: if any consequential contrast predicts repair equally
well, the specifically public-update attribution fails.

\section{Held-out tests and defeat conditions}
\label{sec:identifying-operators}

The theoretical answer generates three projections. \term{Selective fault
attribution} is primary: the independently assigned grammatical-choice profile
should predict whether held-out trouble is directed at grammar. \term{Moving-edge
attribution} predicts greater variation and migration among grammar, lexicon,
and style where the source profile is intermediate or changing. The
\term{operator-profile hypothesis} is downstream: among eligible paradigms,
public-update profile should predict the repair difference between wrong values
and value-preserving realization errors. The first two address the question in
the title; the third tests one proposed consequence within grammar.

Freezing the source profiles isn't enough. The outcome codebook also has to be
fixed before responses are inspected. Prospective elicitation and corpora with
disjoint profiling and outcome partitions can implement the same inferential
order.

\textsc{Study A, Stage 1: choice profiling.} Candidate dependencies are sampled
without first sorting them into grammatical and non-grammatical bins. For each
declared community, situation, medium, and period, production and descriptive
evidence specify the expression--value relation, constructional site,
contrastive frame, alternatives, selection and omission restrictions,
licensing chain, productivity, and population distribution. Negative labels,
corrections, repair, and the held-out items are withheld. Profile components
remain visible; a scalar or ordinal summary, weighting rule, missing-data rule,
and uncertainty treatment are fixed for the particular study rather than
declared universal.

Published descriptions count only through their underlying positive evidence:
attested production, substitution and omission contrasts, conventional
interpretation, productivity, and distribution across independently specified
opportunities. A grammarian's asterisk or category label is attribution evidence
and can't be recycled as source-profile evidence.

The comparison set should include strong paradigms and constructions,
value-preserving realization dependencies such as past-tense exponence,
item-restricted constructional chains, open lexical choices, phonological and
orthographic relations, indexical or stylistic pairings, and social-act cases.
The deviations have to remain interpretable. Frequency, perceptual distance,
processing load, task framing, schooling, and sanction are restricted, matched,
or represented as rival predictors. Boye--Harder diagnostics and
Abrus\'an-inspired semantic repairability are measured only in the domains where
their theories apply.

\textsc{Study A, Stage 2: frozen attribution.} New items, speakers, or a disjoint
corpus partition are first coded for the apparent target: assembly relation,
selected content, expression shape, spelling, situational or indexical fit,
social act, mixed target, or unresolved trouble. A construction-directed
response isn't automatically grammar-directed, because lexical constructions
are part of the boundary to be explained. Open descriptions or repairs should
precede forced labels so that a classroom vocabulary isn't built into the
response. Explicit use of \enquote{grammar} is a separate, norm-sensitive
outcome. Compound responses remain compound, and reliability and uncertainty
are reported.

The primary test asks whether the frozen choice profile improves held-out
spontaneous assembly-directed attribution, and separately explicit grammar
labels, over the declared rivals by at least the prospectively chosen useful
margin. The moving-edge test asks whether
independently intermediate or changing profiles predict greater attributional
variation and the declared direction of movement. A miss defeats or splits the
corresponding claim. In particular, sharp school-taught labels with weak source
profiles support institutional transmission as a cause of the read-out, not a
post-hoc upgrade of the relation.

\textsc{Study B, Stage 1a: strict-paradigm identification.} For the narrower
operator comparison, descriptive and corpus evidence first specifies the
conditioned repertoire, recurring domain, alternatives, and semantic frame. It
then tests whether the set is closed and whether a value has to be selected
whenever the domain applies. Sharp judgments, correction, and the later repair
task are withheld. Tense--aspect systems, agreement paradigms,
switch-reference markers, interrogative particles, and complementizer
inventories are candidates, not presumed cases. This triangulation implements
Poulsen's caution: widespread production is strong evidence for a construction,
but nonoccurrence, rare production, processing cost, and acceptability each
underdetermine stable conventional knowledge when considered alone
\citep[15--20]{poulsen2012grammaticality}.

Closure requires a bounded value inventory that exhausts the choices found in the specified domain across additional corpus samples and, where needed, consultant elicitation. Productive new exponents can realize existing values; a productive route to new values, or an open residual \enquote{other} response, defeats closure. The conditioned repertoire and stopping rule are fixed before the outcome task.

Obligatory selection requires every independently accepted token in the domain to receive one value from that inventory. A zero realization counts only where its distribution contrasts with overt members in the same frame and independently changes the frame interpretation. Mere lack of an overt form doesn't establish a zero value; it may instead show that the proposed domain isn't active.

Dropping N{\o}rg{\aa}rd-S{\o}rensen and Heltoft's markedness criterion enlarges the possible comparison set. Each study reports the primary four-criterion analysis and a sensitivity analysis restricted to paradigms that also show their marked--unmarked asymmetry. A result that disappears in the strict subset supports only the relaxed comparison and can't be attributed to their full architecture.

\textsc{Study B, Stage 1b: operator profiling.} For each identified paradigm, preregister the update variables its values might affect: move type, commitment holder or source, discourse-participant or referent continuity, and response relevance. A separate sample or corpus partition supplies minimal value contrasts while holding at-issue content as constant as the language permits. Classification, interpretation, or independently coded distributional measures estimate how much changing the value changes each component.

Stage~1b yields a four-component vector. A concrete study may preregister an
ordinal comparison or a scalar summary, but the theory supplies no universal
weights. Zero effects remain visible rather than being dropped as
\enquote{inapplicable}; indicator coding, scaling, missing-data rules, contexts,
and uncertainty aggregation are frozen before repair outcomes are inspected.

Profiling evidence stops at the target utterance. It can include blinded coding
of its conventional move type, commitment source, reference continuity, and
licensed response potential, or comprehension probes that classify those
properties. It can't include negative judgments, trouble reports, correction
wording, repair initiation, or the observed next turn. If Stage~2 uses
interactional outcomes, disjoint next-turn samples aren't independent profiling
evidence. This exclusion is stricter than holding out items from the same task.

This step has an information-theoretic interpretation only after the linguistic
codebook has been specified. Entropy reduction or divergence among response
options judged compatible with an utterance can quantify how strongly a value
constrains possible updates
\citep{shannon1948,coverthomas2006}; it can't determine whether the alternatives
form a grammatical paradigm or whether a departure is normatively classified as
an error. Information theory measures a consequence of the supplied value
contrast. It doesn't supply system membership, fault locus, or social
accountability, and token surprisal remains a separate covariate.

Study B Stage~2 uses a different task format and new items, such as open next-turn
responses, production repair, or a held-out interactional corpus partition. Two
theoretical rivals are measured in Study B Stage~1. The first is a semantic-repairability
index motivated by Abrus\'an and colleagues: blinded coders assess whether
at-issue content, predicate--argument structure, and event identity can be
recovered through a local meaning shift, together with the locality and
independently diagnosed type level of the conflict, without the four update
labels or a repair prompt \citep{abrusan2019,abrusanetal2021shift}. The second is
Boye and Harder's discursive-secondary-status diagnostics. The component scores
of semantic repairability remain visible rather than being hidden in a single
post-hoc scale. The analysis reports both rivals' correlations with the operator
profile before fitting the confirmatory model. If the design can't distinguish
these predictors at the achieved support, the study is non-identifying: it
confirms neither the operator attribution nor its failure. Broad scope and
interpersonal consequence are insufficient on their own.

Abrus\'an's measure is N/A, not zero, where no independent semantic type or
composition conflict exists. Pure realization and ordinary licensing failures
thereby fall outside that theory's scope. Boye--Harder diagnostics directly
address grammatical status; their use as a predictor of next-turn repair is an
explicit extension rather than a test of their account as a whole.

Sharing conceptual dimensions doesn't make the outcome definitional. Stage~1b
classifies what a value conventionally configures; Stage~2 asks whether a
respondent directs trouble there. A high-profile wrong value can be accommodated,
ignored, sanctioned, or met with form-only correction. Such responses remain in
the denominator and count against the predicted repair-locus difference.

\textsc{Design covariates.} Opportunity is the number of occasions on which a paradigm value and its competitors could plausibly have been selected. Exponent frequency, surprisal, constructional frame, phonetic salience, processing load, schooling or norm exposure, and social sanction are measured separately. Wherever possible, restriction and matching handle these item- or token-level variables before estimation. The confirmatory model isn't a regression on the whole list: it contains the frozen profile, error locus, their interaction, the two pre-specified theoretical rivals at their appropriate levels, and only adjustments retained by a prospective design simulation. These are design commitments, not controls the theoretical paper claims already to have achieved. Opportunity predicts the concentration of licensing evidence, not paradigm status or operator profile \citep{reynolds2026ovmg}.

Correction is socially organized, so response cause can't be inferred from
response wording. Activity, medium, register, participant role, learner status,
schooling, and the norm centre made relevant in the interaction are specified
before sampling \citep{cameron2012verbalhygiene}. A form-focused correction can
be caused by institutional authority; an update-directed repair can arise from
ordinary mishearing. The outcome code locates the apparent target, while the
design variables address causes. Mixing incompatible norm centres changes the
population target rather than merely adding noise.

\textsc{Study B, Stage 2: frozen outcomes.} Before coding held-out responses, define \term{update confusion}, \term{update-directed repair}, \term{form correction}, \term{constructional correction}, \term{social sanction}, and \term{unresolved trouble}. Update-directed repair targets move type, commitment, discourse-participant or referent continuity, or response relevance; form correction supplies a licensed realization and preserves the intended update. Double-code ambiguous responses and report reliability rather than forcing every case into one bin. The confirmatory comparison includes only paradigms for which independent evidence supports both a genuine value mismatch and a plausible value-preserving realization mismatch; the design can't manufacture the latter merely to retain a paradigm. Eligibility and its relation to expressive substance are reported before profiling. The confirmatory claim concerns this eligible population, not paradigms for which a value-preserving control is unavailable. Those paradigms require a separately labelled extension with a different form or constructional control.

One cross-paradigm interaction supplies the sole confirmatory test for Study B: does the wrong-value versus value-preserving-realization difference in update-directed repair increase with the frozen operator profile after the pre-specified restrictions, matches, and limited adjustments? The within-paradigm difference is a manipulation check, repair locus is the primary outcome, and held-out comprehension is secondary. Repeated responses don't substitute for independent paradigms when estimating this interaction. Before calling a study confirmatory, design simulation has to show that the planned paradigm sample can distinguish the minimum effect from zero at the declared uncertainty threshold. An underpowered or non-identifying corpus analysis is a pilot and licenses neither confirmation nor failure. In an adequately powered identifying study, an absent or smaller interaction defeats the operator-profile hypothesis even if paradigm descriptions and operator profiles remain useful.

\textsc{Projective declarations and revision rules.} Study A projects from a
frozen choice profile to grammar-directed attribution for held-out candidate
dependencies in the declared community, situation, medium, and period. Its
moving-edge component projects from independently intermediate or changing
profiles to attributional variation or migration across declared groups or
times. Study B projects from the frozen profile of an eligible paradigm to its
held-out wrong-value--realization difference in update-directed repair. Each
study declares a minimum useful effect and uncertainty threshold. The
prospective comparisons, rival measures, and separation of tasks supply the
warrant; they aren't part of the world-side profiles. Failure retires or splits
the corresponding attribution for that scope. A preregistered narrower analysis
may be reported as a different claim, but it doesn't rescue the broader one
\citep{Goodman1955,reynolds2026notEveryStableCluster}.

No newly recruited participant sample is logically required for an initial
pilot. Existing correction, interactional, learner, metalinguistic, and usage
data can be divided into disjoint profiling and outcome partitions,
supplemented by descriptive grammars and independent distributional evidence.
Such corpora are selected records of when people correct, but they aren't clean
samples of community classification. Naturally occurring wrong-value and
realization errors may also be too sparse or unevenly distributed for
identification. A corpus study can validate search, profile placement,
eligibility, and coding procedures without being promoted into a test the data
can't support.

\section{Predictions and research strategies}
\label{sec:predictions}

The paper makes three prospective claims, all conditional on the pre-outcome
identification and coding stages in \S\ref{sec:identifying-operators}. The first
two answer the selective-classification question; the third isolates the
operator-specific consequence.

\textsc{(1) Selective fault attribution.} Independently stronger
grammatical-choice profiles should produce more consistent grammar-directed
attribution for held-out interpretable deviations. The comparison isn't
\enquote{grammar versus random bad sentences}. It has to include licensed but
poor lexical choices, phonological and orthographic deviations, situationally
misplaced indexical values, social-act violations, constructional failures,
wrong paradigm values, and value-preserving realization errors. The prediction
fails if frequency, processing, schooling, sanction, semantic repairability, or
discursive secondary status supplies the useful allocation and the choice
profile doesn't.

\textsc{(2) Moving-edge attribution.} Independently intermediate and changing
profiles should produce greater variability and directed migration among
grammar, lexicon, and style attributions. An item-specific restriction embedded
in a constructional chain should be less consistently grammar-directed than a
schematic obligatory dependency but more so than an open collocation, other
conditions matched. Across a declared period, a weakening dependency should
lose spontaneous grammar-directed attribution as recurrence, embedding, or
licensing concentration declines. Persistent classroom labels may lag; that
lag supports institutional transmission as a response cause, not continued
strength of the source profile.

\textsc{(3) The operator-profile-by-error-locus interaction.} Conversation favours self-repair over other-correction: other-correction is dispreferred, while other-initiation of repair is common and usually gives the original speaker an opportunity to repair \citep{schegloffjeffersonsacks1977,dingemanseetal2015repair}. Within each eligible identified paradigm, wrong-value items should disproportionately produce clarification of force, commitment, role, or reference; value-preserving realization errors should disproportionately produce form correction with the intended update retained. That within-paradigm difference checks the manipulation. The confirmatory coefficient is its interaction with the independently assigned operator profile across eligible paradigms.

A higher correction rate for high-profile paradigms is too weak a prediction. *\mention{Depend of} $+$ \textsc{obj} can elicit \enquote{you mean X} through constructional licensing, and a socially marked form can elicit person-targeted sanction. The confirmatory outcome is the locus of the response, not correction rate alone. Cross-linguistic repair research supplies broader coding precedents but not validation of the proposed interaction; Fox and colleagues specifically code same-turn syntactic self-repair rather than the held-out next-turn categories used here \citep{foxetal1996repair,dingemanseblythedirksmeyer2014,dingemanseetal2015repair}. Adult interaction with little formal-school correction can reduce one prescriptive rival, but dialect, learner status, institutional role, and task format remain measured accountability conditions.

Two stress tests concern Claim~3. \textsc{Local redundancy.} In gender, agreement, and reference-tracking systems, a manipulation can leave one token's interpretation unchanged while the paradigm remains informative across its opportunities. After Study B Stage~1b has assigned the profile, held-out contexts should show token-level null effects where lexical semantics already determines the referent and positive effects where reference is underdetermined. The local null is predicted; reassigning the profile after seeing it isn't permitted. A low-profile paradigm that produces the same value-versus-realization repair difference as a high-profile paradigm counts against the hypothesis.

\textsc{Expressive substance.} Where tone, intonation, or nonmanual marking realizes values in an identified paradigm and both mismatch types are independently supported, the same profile-by-error-locus interaction should appear after opportunity, frequency, and perceptual salience are matched \citep{hyman2011,ladd2008}. \textcite{tsiwahetal2021} report a P600-like response to one temporal-agreement violation in Akan, where time reference is expressed through grammatical tone; this prediction concerns the broader matched repair profile, not a unique ERP component. Compare tonal tense mismatches, segmental tense or agreement mismatches, and sociophonetic tone shifts. Wrong values may prompt reanalysis of material in scope; value-preserving realization errors should concentrate correction on the form. This interaction is more specific than the familiar contrast between morphosyntactic and conceptual anomaly \citep{kutashillyard1980,osterhoutholcomb1992,friederici2011}.

The failure conditions belong to the contribution statement. The general answer
fails if the frozen choice profile adds no useful held-out prediction of fault
attribution. The moving-edge claim fails if independently ordered or changing
profiles don't predict attributional variation or direction. In an adequately
powered identifying study, the operator-profile hypothesis adds no explanatory
value if the frozen profile doesn't predict the held-out
value-versus-realization repair difference after the pre-specified design
restrictions, semantic repairability, and discursive secondary status. Source
descriptions and component profiles may remain useful; the failed projection
doesn't. Inseparable rivals or insufficient independent cases make a study
inconclusive, not supportive by default.

The same inferential order applies whether a test uses new elicitation or
disjoint partitions of existing corpora: profile candidate dependencies; where
relevant, identify strict paradigms and public-update effects; freeze rivals and
outcome codes; and only then inspect the held-out attributions or repairs.

\section{Conclusion}

Why, then, does one interpretable deviation count as ungrammatical while another
doesn't? The proposed answer is that speakers imperfectly track the relation
that failed. Grammar consists centrally of conventional form--meaning pairings,
but neither conventionality, pairing, nor systemhood is sufficient. The
grammatical core consists of relations used recurrently to assemble expressions:
constructions, relational values, obligatory choices, and their licensed
realizations. Strict paradigms anchor that core; construction-sized and
item-restricted dependencies produce layered cases. Selecting a poor but
licensed word, style, accent, or social act can instead preserve the relevant
assembly and make content, situational fit, persona, or conduct the fault target.

That's a graded, scoped answer rather than an exceptionless definition. A
candidate relation receives a profile from form--value coupling, recurrence,
contrastive choice architecture, and conditioned community licensing without
using the labels to be predicted. A stronger
profile is then expected to support more consistent grammar-directed attribution
in held-out cases. The worldly profile, our uncertainty about it, and the
sharpness of a respondent's label remain distinct. Edges can also move with
community, activity, medium, and time.

The motivating French example makes the point precisely. In ordinary connected
Reference French, *\olang{le hiver} preserves referent, definiteness, number, and
gender while using an unelided realization before an \term{h muet} noun. The
dependency is morphophonological, but it's embedded in article realization.
That makes grammatical-form attribution plausible without turning the item into
a concord or wrong-value case. It also makes the example a genuine interface:
speech, writing, pausing, learner correction, and norm centre can shift
attribution toward phonology, orthography, or grammatical realization. Those
possibilities test the profile account rather than disappearing by definition.

OVMG contributes the token-level and epistemic decomposition. Coverage,
compatibility, saturation, and licensing locate distinct possible failures;
population status, evidence, and subjective read-out aren't one variable.
Opportunity and strong competitors can explain why a transparent
value-preserving realization becomes sharply excluded. They don't decide the
fault label from frequency alone, and institutional teaching or stigma can make
the label mis-track the source profile.

Public-update contribution adds a narrower claim. Among independently identified
paradigms, values differ in how strongly they configure move type, commitment,
reference continuity, and response relevance. Across paradigms that genuinely
support both error types, a stronger frozen public-update profile is predicted
to increase the difference in update-directed repair between wrong values and
value-preserving realization errors. That projection explains neither
*\olang{le hiver} nor grammar as a whole.

The proposal is cumulative. Abrus\'an and colleagues explain one major boundary
through type level, locality, and semantic shiftability. Coseriu supplies a
longstanding architecture of system, norm, speech, and evaluative domains.
Heltoft shows how paradigmatic grammar organizes illocutionary and
interlocutionary relations, while N{\o}rg{\aa}rd-S{\o}rensen and Heltoft provide
the domain--frame scaffold adopted here. Boye and Harder's discursive secondary
status remains a principled alternative predictor. The residual contribution is
not a new ingredient but a cross-level decomposition and three prospective
comparisons that can distinguish their reach.

The proposal can lose in three ways. If the frozen
grammatical-choice profile doesn't improve held-out fault attribution beyond
semantic repairability, discursive secondary status, frequency, processing,
schooling, and sanction, the general answer fails or has to be split. If
independently intermediate or changing profiles don't predict variable or
moving attribution, the boundary claim fails. If the public-update profile does
not predict the held-out value--realization repair difference in an identifying
study, the operator-specific attribution fails. Descriptive analyses may
survive each result; they may not be used to rename a failed projection as
success.

Clausal architecture remains important because it can host paradigms for clause
type, commitment, scope, and reference tracking. Tone, gesture, particles, word
order, and inflection can realize those values. None of this makes every
constraint on a clause high-profile, and the stronger historical suggestion
that public-update work preferentially recruits paradigms remains unconfirmed.

\section*{Acknowledgements}

Jamie Ramsden's observation about French article realization prompted the argument.

\clearpage
\printbibliography
\end{document}
